\documentclass[a4paper,14pt]{extarticle}

\input{C:/preamble.tex}

\usepackage{tcolorbox}

\usepackage{amsmath}
\usepackage{amssymb}
\usepackage{mathtools}

\begin{document}
	
	\begin{center}
		\textbf{Модели фильтрации крови при наличии тромба}
	\end{center}
	
	В рассмотренных работах течение крови в сосуде с тромбом
	изучается с помощью фильтрационных моделей. Тромб
	представляется проницаемой пористой средой. В отличие от
	задачи сопряжения Стокса--Дарси, движение крови вне тромба
	в этих работах также описывается фильтрационной моделью,
	но с большей проницаемостью.
	
	\section{Стационарная фильтрация и разрушение тромба}
	
	В работе~\cite{stationary} рассматривается течение крови
	в сосуде при наличии прикреплённого к его стенке пористого
	тромба. Сосудистая область и пространство пор тромба
	характеризуются различными проницаемостями. Для описания
	движения используется закон Дарси
	\begin{equation}
		\mathbf q_i=-\frac{k_i}{\mu}\nabla p,
		\qquad i=1,2,
	\end{equation}
	где $p$ --- давление, $\mu$ --- вязкость крови, $k_i$ ---
	проницаемость соответствующей области, $\mathbf q_i$ ---
	скорость фильтрации. В модели также рассматриваются
	уравнения баланса массы и возможное изменение объёмной доли
	материала тромба.
	
	Разрушение материала связывается с величиной градиента
	давления $|\nabla p|$: предполагается, что оно начинается
	при превышении некоторого критического значения $I_0$.
	Для характеристики $N$ разрушаемого материала записано
	уравнение
	\begin{equation}
		\frac{\partial N}{\partial t}
		=-\Psi\bigl(N,|\nabla p|-I_0(N)\bigr),
	\end{equation}
	где функция $\Psi$ задаёт скорость разрушения. В рассматриваемом
	приближении, если изменение $N$ не влияет на фильтрационные
	свойства среды, сначала можно найти поле давления и его
	градиента, а затем рассчитать изменение $N$.
	
	В вычислительных экспериментах исследуется влияние формы
	и проницаемости тромба на давление, скорость фильтрации и
	градиент давления. Согласно результатам авторов, форма
	тромба и угол его примыкания к стенке заметно влияют на
	расположение области максимального градиента давления.
	При остром угле в рассмотренных примерах значительный
	градиент возникает около основания тромба, что связывается
	с риском его отрыва. Проницаемость также влияет на величину
	градиента и перепад давления в сосуде.
	
	\section{Нестационарная фильтрация}
	
	В работе~\cite{transient} рассматривается фильтрация крови
	в сосуде с кольцевым тромбом при периодическом изменении
	давления, связанном с работой сердца. Задача записывается в
	двумерном продольном сечении сосуда. Проницаемость $K(x,z)$
	считается заданной, причём внутри тромба она может быть
	неоднородной и существенно меньше, чем вне его.
	
	Модель включает уравнение баланса и закон Дарси:
	\begin{equation}
		\beta\frac{\partial P}{\partial t}
		+\operatorname{div}\mathbf V=0,
		\qquad
		\mathbf V=-\frac{K(x,z)}{\mu}\nabla P,
	\end{equation}
	где $P$ --- давление, $\mathbf V$ --- скорость фильтрации,
	$\beta$ --- параметр, учитывающий сжимаемость системы, а
	$\mu$ --- вязкость. После подстановки закона Дарси получается
	уравнение для давления:
	\begin{equation}
		\beta\frac{\partial P}{\partial t}
		-\operatorname{div}\left(\frac{K(x,z)}{\mu}\nabla P\right)=0.
	\end{equation}
	
	На входе и выходе задаётся зависящее от времени давление;
	на остальных границах рассматриваются условия отсутствия
	нормального потока. Для решения используется сетка в
	продольном сечении и неявная конечно-разностная схема.
	Авторы сравнивают распределения давления, скорости и
	градиента давления в течение сердечного цикла для различных
	форм и структур тромба.
	
	В приведённых расчётах неоднородность проницаемости
	существенно изменяет распределение градиента давления
	внутри тромба. Его максимум может достигаться в области
	низкой проницаемости; в отдельных рассмотренных примерах
	он более чем в четыре раза превышает максимум для
	однородного тромба. Авторы также отмечают влияние формы
	тромба и периодического изменения давления на возможное
	разрушение материала.
	
	\section{Общее замечание}
	
	В обеих работах градиент давления используется как
	характеристика нагрузки на материал тромба и возможного
	риска его разрушения. Полученные зависимости относятся к
	выбранным моделям и параметрам расчёта. Дальнейшее развитие
	может состоять, например, в изменении закона проницаемости,
	учёте её связи с разрушением тромба или применении другого
	численного метода.
	
	\begin{thebibliography}{9}
		\bibitem{stationary}
		Kramchenkov M.~G., Konyukhov V.~M. и др.
		\textit{Filtration Model of a Blood Flow in the Circulatory
			System at the Risk of Clot Breakage}.
		
		\bibitem{transient}
		Konyukhov I.~V., Konyukhov V.~M. и др.
		\textit{Transient processes of the blood flow in a vessel at
			the presence of circular clot}.
	\end{thebibliography}
	
\end{document}
